\section{结论}
\label{sec:ink-conclusion}

本文提出 \method{}，根据专家与当前合并模型之间尚未消除的预测残差，动态识别每一层需要修正的功能方向。教师--学生 KL 关于层表示的归一化负梯度给出局部修正，其二阶矩与当前输入激活几何共同形成 Kronecker 因子化正规方程，并由共轭梯度、逐块保守步长和留出教师 KL 完成迭代合并。冻结的 CLIP ViT-B/32 八任务实验表明，该方法在两个独立校准种子下分别取得 $86.95\%$ 和 $87.05\%$ 的平均准确率。

当前结论限定于可精确枚举类别、具有每任务无标签校准样本的 CLIP 视觉基准。进一步验证仍需覆盖更多架构、任务规模与随机种子，并评估校准数据需求和迭代统计重估的计算代价。
